%% sections/experiments.tex

\section{Experimental Setup}
\label{sec:experiments}

\subsection{Datasets}

Training and evaluation used two complementary corpora: the official PAN-2026 corpus and the RAID benchmark~\cite{dugan:2024}.

\subsubsection{PAN-2026 Corpus}
The organiser-provided dataset spans three genres—\textit{essays}, \textit{news}, and \textit{fiction}—and covers multiple generative models including GPT-4o, GPT-4o-mini, Claude 3 Sonnet, Llama-3.1, Mistral, and others. Text generation was conditioned on style prompts referencing real human authors to make the AI output harder to detect. The label distribution is slightly imbalanced: 61.7\% AI vs.\ 38.3\% human.

\subsubsection{RAID Dataset}
RAID~\cite{dugan:2024} is a large-scale adversarial benchmark covering 11 domains, 11 generative models, and 12 attack strategies. We extracted text from three domain groups that overlap with the PAN genres: \textit{news}, \textit{books}, and \textit{poetry}. We further restricted the adversarial attack conditions to five types that represent realistic transformations:
\begin{itemize}
  \item \textbf{None}: raw model output without post-processing.
  \item \textbf{Homoglyph}: character substitution with visually similar Unicode characters.
  \item \textbf{Synonym}: lexical substitution using a paraphrase model.
  \item \textbf{Paraphrase}: full-sentence paraphrase.
  \item \textbf{Alternative spelling}: regional or archaic spelling variants.
\end{itemize}

After filtering, we sampled 30,000 records for training (balanced 50/50 human/AI) and 10,000 for testing from the original 5.6 million RAID rows.

\subsubsection{Combined Training Set}
The main training split merges all 23,682 PAN-2026 training records with a 15,000-record sample from the RAID training split, yielding 38,682 records. Combining the two corpora improves cross-domain generalisation, as the PAN data provides style-mimicry examples while RAID provides adversarial attack variety. The validation set combines the full 3,589 PAN-2026 validation records with 14,834 RAID records that were held out of training.

Table~\ref{tab:datasets} summarises the dataset statistics.

\begin{table}[htbp]
  \centering
  \caption{Dataset composition used in our experiments.}
  \label{tab:datasets}
  \setlength{\tabcolsep}{5pt}
  \begin{tabular}{lrrrrr}
    \toprule
    \textbf{Split} & \textbf{Source} & \textbf{Total} & \textbf{Human} & \textbf{AI} & \textbf{AI\%} \\
    \midrule
    Train & PAN-2026 & 23,682 & 9,077 & 14,605 & 61.7 \\
    Train & RAID (sample) & 15,000 & 7,500 & 7,500 & 50.0 \\
    \midrule
    \textit{Train (combined)} & & \textit{38,682} & \textit{16,577} & \textit{22,105} & \textit{57.1} \\
    \midrule
    Val & PAN-2026 & 3,589 & 1,376 & 2,213 & 61.7 \\
    Val & RAID (held-out) & 14,834 & 7,422 & 7,412 & 50.0 \\
    \midrule
    \textit{Val (combined)} & & \textit{18,423} & \textit{8,798} & \textit{9,625} & \textit{52.2} \\
    \midrule
    Test & RAID (held-out) & 10,000 & 4,976 & 5,024 & 50.2 \\
    \bottomrule
  \end{tabular}
\end{table}

\subsection{Data Preprocessing}

All texts undergo a standardised pipeline before training:
\begin{enumerate}
  \item \textbf{Case normalisation}: convert to lowercase.
  \item \textbf{PII masking}: replace email addresses, Twitter usernames, and phone numbers with the tokens \texttt{[EMAIL]}, \texttt{[USER]}, and \texttt{[PHONE]}, respectively.
  \item \textbf{Deduplication}: remove exact-duplicate records (same text after normalisation).
  \item \textbf{Minimum-length filtering}: discard records with fewer than 50 characters.
\end{enumerate}

Stylometric features are extracted from the normalised text prior to further tokenisation. The transformer tokeniser (ModernBERT WordPiece vocabulary) then encodes the text up to a maximum of 512 subword tokens, truncating longer texts.

\subsection{Training Configuration}

Both training stages share the same optimiser settings to allow direct comparison. Hyperparameters were selected based on a preliminary grid search over learning rate $\in \{5\times 10^{-6}, 2\times 10^{-5}, 5\times 10^{-5}\}$ and batch size $\in \{16, 32\}$.

\begin{table}[htbp]
  \centering
  \caption{Training hyperparameters for both stages.}
  \label{tab:hyperparams}
  \begin{tabular}{ll}
    \toprule
    \textbf{Parameter} & \textbf{Value} \\
    \midrule
    Optimiser           & AdamW (fused CUDA kernel) \\
    Learning rate       & $2 \times 10^{-5}$ \\
    LR schedule         & Cosine annealing \\
    Warmup ratio        & 3\% of total steps \\
    Weight decay        & 0.01 \\
    Gradient clipping   & 1.0 \\
    Batch size (train)  & 32 \\
    Batch size (eval)   & 64 \\
    Max sequence length & 512 tokens \\
    Epochs              & 10 \\
    Random seed         & 42 \\
    Loss function       & BCEWithLogitsLoss \\
    Positive class weight ($\omega$) & 0.746 \\
    Mixed precision     & bfloat16 \\
    \bottomrule
  \end{tabular}
\end{table}

The positive-class weight $\omega = 0.746$ was computed from the training label distribution as $N_{\text{neg}} / N_{\text{pos}} = 16{,}577 / 22{,}105$, which upweights the human class to counteract the 57\% AI prevalence.

Stage~2 initialises from the best Stage~1 checkpoint (selected by highest mean score on the combined validation set) and continues with identical hyperparameters with all parameters unfrozen. Mixed-precision training with bfloat16 is used throughout to reduce memory consumption.

\subsection{Ablation Experiments}

In addition to the two-stage full system, we conducted two single-dataset ablation studies to isolate the contribution of each data source:

\begin{itemize}
  \item \textbf{RADAR-PAN26}: trained with frozen encoder on the PAN-2026 corpus alone (no RAID data). This evaluates the system's performance under domain-in-distribution conditions without additional data augmentation.
  \item \textbf{RADAR-RAID}: trained with frozen encoder on the RAID corpus alone (no PAN-2026 data). This isolates the benefit of adversarial-attack coverage.
\end{itemize}

Both ablations follow Stage~1 training (frozen encoder) only, matching the conditions of the full Stage~1 baseline.

\subsection{Hardware and Software}

All experiments were conducted on a single NVIDIA A100-SXM4-40GB GPU. The implementation uses PyTorch~2.0 and the HuggingFace Transformers library~\cite{wolf:2020}. Models were saved as HuggingFace Hub repositories for reproducibility.

\subsection{Evaluation Protocol}

Models are evaluated at the end of each epoch on the combined validation set. The best checkpoint is selected based on the mean of all five PAN metrics. Final evaluation on the RAID held-out test set is performed once at the end of training to report cross-domain generalisation. PAN-2026 official evaluation runs on the organiser's hidden test set through the TIRA~\cite{froebe:2023b} platform. The validation-set metrics reported in the next section serve as approximations of the expected test-set performance.
